\mxnewpage\section*{Correction EV3-B}\label{EV3-B-CORR}
\textbf{QCM :} 1 A, 2 B, 3 A, 4 B, 5 D, 6 B, 7 C, 8 C, 9 C, 10 D, 11 D, 12 D. Chaque réponse correcte : 0{,}5 point ; total 6.
\paragraph{Vérifications des automatismes.} Priorités ; somme de fractions ; soustraction des exposants ; coefficient multiplicateur 1{,}2 ; division par 0{,}8 ; produit 1{,}25 fois 0{,}8 ; résolution affine ; accroissement vertical divisé par l’horizontal ; moyenne ; angle divisé par 360 ; complément à 1 ; substitution entre parenthèses.
\subsection*{Exercice 1 -- Quatre essais et un gain}
\begin{enumerate}
\item $(\dfrac{1}{2})^4=\dfrac{1}{16}$. \textit{Barème : 1 point. Résultat et justification élémentaire attendus ; demi-point possible pour démarche correcte inachevée.}
\item $E(G)=\dfrac{1}{2}$, $E(G^2)=4$ ; $V(G)=\dfrac{15}{4}$. \textit{Barème : 2 points. Méthode ou justification : 1 ; calcul et conclusion : 1.}
\item Mise $\dfrac{1}{2}$. L’espérance du gain brut vaut $\dfrac{1}{2}$, non possible sur un seul jeu. \textit{Barème : 2 points. Méthode ou justification : 1 ; calcul et conclusion : 1.}
\end{enumerate}
\subsection*{Exercice 2 -- Une dérivée à signe constant}
\begin{enumerate}
\item $f\prime=(x^2+2x+1)e^x=(x+1)^2e^x$. \textit{Barème : 1 point. Résultat et justification élémentaire attendus ; demi-point possible pour démarche correcte inachevée.}
\item Positive sauf en -1 ; f croît strictement (sur les deux côtés et par raccord). Pas d’extremum en -1. \textit{Barème : 2 points. Méthode ou justification : 1 ; calcul et conclusion : 1.}
\item $f(0)=1,f\prime(0)=1$ : $y=x+1$ ; $f(1)=2e$. \textit{Barème : 2 points. Méthode ou justification : 1 ; calcul et conclusion : 1.}
\end{enumerate}
\subsection*{Exercice 3 -- Projection sur une diagonale}
\begin{enumerate}
\item Normale (1;-1). \textit{Barème : 1 point. Résultat et justification élémentaire attendus ; demi-point possible pour démarche correcte inachevée.}
\item H(1;1), PH=$\sqrt2$. \textit{Barème : 1 point. Résultat et justification élémentaire attendus ; demi-point possible pour démarche correcte inachevée.}
\item Angle $\dfrac{\pi}{4}$ ; vecteur $(\dfrac{\sqrt 2}{2};\dfrac{\sqrt 2}{2})$. \textit{Barème : 2 points. Méthode ou justification : 1 ; calcul et conclusion : 1.}
\end{enumerate}
\paragraph{Étalonnage.} Durée cible : 20–25 min pour les automatismes, 80–85 min pour les exercices, 10–15 min de relecture. Ces durées sont des estimations éditoriales à mesurer en classe. Accepter toute méthode correcte ; ne pas sanctionner deux fois une erreur numérique propagée avec cohérence. Le barème proposé ne vaut pas décision d’harmonisation de l’équipe.
